\chapter{भाग}\label{ch:g4:division}

पिछले साल हमने पहाड़ों और छोटे शेषफलों के साथ भाग दिया था
(\cref{ch:g3:division}); यह अध्याय पूरा लिखित तरीक़ा सिखाता
है --- लंबा भाग --- जो किसी भी संख्या को एक-एक अंक करके सँभाल
लेता है।

\section{लंबे भाग का तरीक़ा}

\begin{method}[एक अंक की संख्या से लंबा भाग]\label{met:g4:division:method}
$749$ में $6$ का भाग देने के लिए:
\begin{enumerate}
  \item सबसे बाईं ओर का अंक लो, $7$ (सैकड़ा): $7$ में $6$ एक
        बार जाता है; \omterm{def:g3:division:remainder}{भागफल} में $1$ लिखो, $6$ घटाओ,
        \omterm{def:g3:division:remainder}{शेषफल} $1$;
  \item अगला अंक, $4$, \emph{नीचे उतारो}: अब $14$ (दहाई) में
        भाग दो: $6 \times 2 = 12$; $2$ लिखो, \omterm{def:g3:division:remainder}{शेषफल} $2$;
  \item $9$ नीचे उतारो: $29$ में भाग दो: $6 \times 4 = 24$;
        $4$ लिखो, \omterm{def:g3:division:remainder}{शेषफल} $5$;
  \item अब कोई अंक नहीं बचा: \omterm{def:g3:division:remainder}{भागफल} $124$ है, \omterm{def:g3:division:remainder}{शेषफल} $5$।
\end{enumerate}
जाँच: $6 \times 124 + 5 = 744 + 5 = 749$, और $5 < 6$।
\end{method}

\begin{omfigure}
\begin{tikzpicture}[scale=0.95]
  % अंकों के स्तंभ x = 0.2, 0.8, 1.4 पर; हर अंक अपने स्थान के नीचे
  \node[font=\normalsize] at (0.2,3.0) {$7$};
  \node[font=\normalsize] at (0.8,3.0) {$4$};
  \node[font=\normalsize] at (1.4,3.0) {$9$};
  % भाजक का कोष्ठक और भागफल
  \draw[thick] (1.8,2.0) -- (1.8,3.35);
  \draw[thick] (1.8,2.62) -- (3.7,2.62);
  \node[font=\normalsize] at (2.35,3.0) {$6$};
  \node[font=\normalsize, omDef] at (2.35,2.2) {$1$};
  \node[font=\normalsize, omDef] at (2.95,2.2) {$2$};
  \node[font=\normalsize, omDef] at (3.55,2.2) {$4$};
  % कदम 1: 7 - 6 = 1
  \node[font=\normalsize, omThm] at (-0.35,2.4) {$-$};
  \node[font=\normalsize, omThm] at (0.2,2.4) {$6$};
  \draw (-0.55,2.1) -- (0.55,2.1);
  \node[font=\normalsize] at (0.2,1.8) {$1$};
  \node[font=\normalsize] at (0.8,1.8) {$4$};
  \draw[-Stealth, omExo, thin] (0.8,2.75) -- (0.8,2.05);
  % कदम 2: 14 - 12 = 2
  \node[font=\normalsize, omThm] at (-0.35,1.2) {$-$};
  \node[font=\normalsize, omThm] at (0.2,1.2) {$1$};
  \node[font=\normalsize, omThm] at (0.8,1.2) {$2$};
  \draw (-0.55,0.9) -- (1.15,0.9);
  \node[font=\normalsize] at (0.8,0.6) {$2$};
  \node[font=\normalsize] at (1.4,0.6) {$9$};
  \draw[-Stealth, omExo, thin] (1.4,2.75) -- (1.4,0.85);
  % कदम 3: 29 - 24 = 5
  \node[font=\normalsize, omThm] at (0.25,0.0) {$-$};
  \node[font=\normalsize, omThm] at (0.8,0.0) {$2$};
  \node[font=\normalsize, omThm] at (1.4,0.0) {$4$};
  \draw (0.05,-0.3) -- (1.75,-0.3);
  \node[font=\normalsize] at (1.4,-0.6) {$5$};
  % टिप्पणी
  \node[right, font=\small, omExo, align=left] at (4.5,1.4)
    {$7 - 6 = 1$, $4$ नीचे उतारो;\\
     $14 - 12 = 2$, $9$ नीचे उतारो;\\
     $29 - 24 = 5$: रुको।\\[3pt]
     भागफल $124$, शेषफल $5$\\
     जाँच: $6 \times 124 + 5 = 749$};
\end{tikzpicture}

{\small $749$ में $6$ का लंबा भाग, पूरा लिखा हुआ: हर घटाव दिखता
है, हर अंक अपने स्तंभ में रहता है, और स्लेटी तीर नीचे उतारे जा
रहे अंक दिखाते हैं।}
\end{omfigure}

\begin{example}[भागफल में शून्य]\label{ex:g4:division:zero}
$618$ में $3$ का भाग दो: सैकड़ा, $6 \div 3 = 2$; दहाई, $1$ नीचे
उतारो: $1$ में $3$ \omterm{def:g1:counting:zero}{शून्य} बार जाता है --- \emph{वह $0$ लिखो}! ---
\omterm{def:g3:division:remainder}{शेषफल} $1$; इकाई, $8$ नीचे उतारो: $18 \div 3 = 6$। \omterm{def:g3:division:remainder}{भागफल}
$206$, \omterm{def:g3:division:remainder}{शेषफल} $0$। बीच का \omterm{def:g1:counting:zero}{शून्य} भूल जाना ($26$ लिखना) आम ग़लती
है; जाँच
$3 \times 26 = 78 \neq 618$ उसे तुरंत पकड़ लेती है।
\end{example}

\begin{example}[भागफल का आकार आँकना]\label{ex:g4:division:size}
$749$ में $6$ का भाग देने से पहले उत्तर को घेर लो:
$6 \times 100 = 600$ और $6 \times 200 = 1\,200$, इसलिए \omterm{def:g3:division:remainder}{भागफल}
$100$ और $200$ के बीच होगा --- वह तीन अंकों का होगा। एक पंक्ति
का यह अंदाज़ा अंक ग़लत जगह रखने की ज़्यादातर ग़लतियाँ रोक देता है।
\end{example}

\section{सवाल क्या माँगता है, वही चुनो}

\begin{method}[भागफल, शेषफल, या भागफल और एक]\label{met:g4:division:interpret}
जैसे \cref{met:g3:division:interpret} में, कहानी ही फ़ैसला करती है:
\begin{enumerate}
  \item ``कितने पूरे \dots'' या ``हर एक को कितने'': \omterm{def:g3:division:remainder}{भागफल};
  \item ``कितने बचे'': \omterm{def:g3:division:remainder}{शेषफल};
  \item ``सबको रखने के लिए कितने पात्र चाहिए'': \omterm{def:g3:division:remainder}{भागफल}, और
        \omterm{def:g3:division:remainder}{शेषफल} \omterm{def:g1:counting:zero}{शून्य} न हो तो उसमें एक और।
\end{enumerate}
\end{method}

\begin{example}\label{ex:g4:division:problem}
$530$ किताबें $8$-$8$ के डिब्बों में रखनी हैं।
\[
530 = 8 \times 66 + 2 .
\]
कितने \emph{पूरे} डिब्बे? $66$। कितनी किताबें पूरे डिब्बे में
नहीं हैं? $2$। \emph{सारी} किताबें रखने के लिए कितने डिब्बे?
$67$। तीन सवाल, एक ही भाग।
\end{example}

\begin{example}[पैसों का बराबर बँटवारा]\label{ex:g4:division:money}
चार दोस्त $92$ के उपहार की क़ीमत बराबर बाँटते हैं:
$92 \div 4 = 23$ ठीक-ठीक (\omterm{def:g3:division:remainder}{शेषफल} $0$)। हर एक $23$ देता है।
\omterm{def:g3:division:remainder}{शेषफल} \omterm{def:g1:counting:zero}{शून्य} हो तो भाग ``पूरा बँट जाता है'' और बँटवारा बिलकुल
बराबर होता है।
\end{example}

\section{अभ्यास}

\begin{exercise}[$\star$]\label{exo:g4:division:1}
हर \omterm{def:g3:division:remainder}{भागफल} को \cref{ex:g4:division:size} की तरह घेरो (किन दो
``गोल'' संख्याओं के बीच?), फिर लंबा भाग करो:
$96 \div 4$; \quad $87 \div 5$।
\end{exercise}

\begin{exercise}[$\star$]\label{exo:g4:division:2}
लंबे भाग करो और हर एक की जाँच करो:
$672 \div 4$; \quad $925 \div 7$; \quad $804 \div 6$।
\end{exercise}

\begin{exercise}[$\star$]\label{exo:g4:division:3}
\omterm{def:g1:counting:zero}{शून्य} से सावधान (देखो \cref{ex:g4:division:zero}):
$816 \div 4$; \quad $420 \div 7$; \quad $2\,512 \div 5$।
\end{exercise}

\begin{exercise}[$\star$]\label{exo:g4:division:4}
जाँच का संबंध ($a = b \times q + r$) लिखो: $85 \div 9$;
\quad $1\,000 \div 3$।
\end{exercise}

\begin{exercise}[$\star$]\label{exo:g4:division:5}
$156$ अंडे $6$-$6$ के डिब्बों में रखे जाते हैं। कितने डिब्बे भरते हैं?
\end{exercise}

\begin{exercise}[$\star$]\label{exo:g4:division:6}
$250$ cm का रिबन $8$-$8$ cm के टुकड़ों में काटा जाता है। कितने
टुकड़े, और कितनी लंबाई बच जाती है?
\end{exercise}

\begin{exercise}[$\star$]\label{exo:g4:division:7}
$375$ विद्यार्थी $52$-$52$ सीटों वाली बसों से खेल देखने जाते
हैं। कितनी बसें चाहिए? (यह
\cref{met:g4:division:interpret} का कौन-सा मामला है? गुणजों की
मदद से $52$ का भाग दो: $52 \times 7 = 364$।)
\end{exercise}

\begin{exercise}[$\star$]\label{exo:g4:division:8}
तीन दोस्त $234$ कंचे बराबर बाँटते हैं। हर एक को कितने कंचे?
गुणा से जाँचो।
\end{exercise}

\begin{exercise}[$\star$]\label{exo:g4:division:9}
भाज्य ढूँढ़ो: $7$ से भाग देने पर \omterm{def:g3:division:remainder}{भागफल} $58$ और
\omterm{def:g3:division:remainder}{शेषफल} $3$ आता है।
\end{exercise}

\begin{exercise}[$\star\star$]\label{exo:g4:division:10}
एक पुस्तकालयाध्यक्ष $438$ किताबें रैक पर लगाता है, हर ख़ाने में
$9$। ख़ाने $6$-$6$ की अलमारियों में आते हैं। कितने ख़ाने चाहिए?
कितनी अलमारियाँ ख़रीदनी होंगी? (दो भाग, दोनों में
``और एक'' वाला सवाल।)
\end{exercise}

\begin{exercise}[$\star\star$]\label{exo:g4:division:11}
$8$ से भाग देने पर \omterm{def:g3:division:remainder}{भागफल} \omterm{def:g3:division:remainder}{शेषफल} के बराबर आता है। संभव भाज्य
कौन-कौन-से हैं? सब गिनाओ। (\omterm{def:g3:division:remainder}{शेषफल} $8$ से नीचे रहना
चाहिए।)
\end{exercise}
